\section{Pet Peeves}

\paragraph{Reputations}


A friend's wealthy boyfriend claimed to know a ``computer guy'' who, he claimed is a ``genius. He sent the genius to set up her vanity email account on her cell phone. After more than an hour, he failed to do so; he left her without access to email. I then had to go over to set up the email correctly. Now I understand it does not take a genius to setup an email account, but the story points out how reputations are made, usually undeservedly.

In another all too brief online discussion, I referred to the first chapter of \emph{Revolt Against the Modern World\footnote{\url{https://gornahoor.net/?p=8408}}}. To my surprise, one fellow, who is internationally recognized as an Evola ``expert'', asked why he should consider Julius Evola as an authority on anything. In that chapter, Evola was describing the foundation of Tradition itself, not asserting his own opinion.

\paragraph{Speaking from the Heart}


Do not be awed by reputations, which are quite often undeserved. Learn to distinguish those who speak from the heart with gnosis rather than from the head with learned erudition.

The former is in the vertical direction: mystical experience leads to direct gnosis which is then, and only then, expressed in words. The word is a revelation of the mystical experience.

Erudition works in the horizontal direction, relying on intellectual studies, and knows about things, i.e., indirect knowledge. It does have its place, but not the first place. Intellectual concepts cannot lead to mystical experience.

\paragraph{Degrees of Education}


Although many prefer to remain in relative ignorance --- which is the natural state of man --- there is the possibility of knowledge. Once the seven liberal arts have been mastered, there are three advanced degrees awarded in specialized disciplines.

These are the three degrees of education and what they should represent:

\begin{itemize}


\item The \textbf{bachelor} learns about all the various schools of his discipline.

\item The \textbf{master} becomes an expert in one of them.

\item The \textbf{doctor} creates his own school.
\end{itemize}


Doctors are rare, but many have the potential to become masters. Nevertheless, most are content to remain bachelors. Bachelors prefer to ``refute'' the work of a great philosopher, for example, with a pithy sentence or two. Or they will characterize epochs lasting centuries, or vast civilizations, in the same way. In short, their understanding is limited to what can fit onto a few placards.

\paragraph{Rhetoric}


Rhetoric is the art of persuasion. There are three modes to persuade an audience:

\begin{itemize}


\item \textbf{Logos}: by the use of logical arguments


\item \textbf{Pathos}: by appealing to emotions


\item \textbf{Ethos}: by appealing to ``the guiding beliefs or ideals that characterize a community, nation, or ideology.''
\end{itemize}


Logos is seldom effective except among mathematicians. As an exercise, the next time you watch a ``news'' show or listen to a politician, take a step back to pay attention to the rhetorical techniques. Some examples follow, without exhausting all the techniques. I will use commonly heard arguments about illegal immigration.

For Logos, they may try to point out the logical inconsistencies of their opponents. Or hardliners will appeal to the positive law, characterizing situations by their legality.

For Pathos, they will exaggerate situations to tug at emotions. Examples are the claim that Trump is ``separating families'', one side calling it cruel and inhumane, with the other asserting its lawfulness. This is very common.

For Ethos, there is the appeal to allegedly moral values. Politicians will cynically bring up ``Christian ethics'', even if they are hardly devout Christians, while, at other times, denying that the USA is a Christian nation. Or to be succinct, they will claim that a policy ``does not represent who we are'' --- as though the USA has a monosemous identity.

\paragraph{God's Self-Sacrifice}


Another commonly heard claim is that the ``church'' declined with the introduction of Cabbalistic methods during the Renaissance. This is another placard, which ignores complexities. As a point of fact, for example, the Orthodox Christian view on creation is virtually identical to the Cabbala.

\textbf{Boris Mouravieff} informs us:


\begin{quotex}
Orthodox Tradition teaches that the Universe was created by a sacrifice of God. We shall understand this postulate better if we consider that it differentiates between the state of manifested Divinity and that of unmanifested Divinity --- which is therefore limitless and free from all conditions. God's sacrifice is Self-limitation by manifestation.
\end{quotex}


This is identical to the Cabbala teaching of \emph{tsimtsum}, or God's withdrawal and self-limitation. \textbf{Valentin Tomberg} explains why:

\begin{quotex}
In order to create the world ex nihilo, God had first to bring the void itself into existence. He had to withdraw within in order to create a mystical space, a space without his presence --- the void. And it is in thinking this thought that we assist at the birth of freedom. The void --- the mystical space from which God withdrew himself through his act of \emph{tsimtsum} --- is the place of origin of freedom, i.e. the place of the origin of an ``ex-istence'' which is absolute potentiality, not in any way determined.
\end{quotex}


Otherwise God would be responsible for sin and evil:

\begin{quotex}
The idea of \emph{tsimtsum}, the withdrawal of God in order to create freedom, and that of divine crucifixion on account of freedom, are in complete accordance. For the withdrawal of God in order to make a space for freedom and his renunciation of the use of his power against the abuse of freedom (within determined limits) are only two aspects of the same idea.
\end{quotex}


God created Eden, not the man-made world of our experience. We know what to do to make it better, yet we choose not to.

\paragraph{Intelligent Children}


I live in a town where all children seem to be above average. Mothers are thrilled at their children's academic success, and are justly proud of their sons' intelligence. Women, if you can believe dating site profiles, are seeking intelligent men to be with.

The irony is that when those sons grow up, women no longer value their intelligence, or more precisely, they value intelligence only up to a certain limit. They certainly don't much value the opinions and ideas of men whose intelligence passes that threshold.

\paragraph{Fool Me Once}


I recently witnessed a discussion about a woman who was replaced by a trophy wife. Hurt, she turned to online dating where she met a man allegedly living far away, although never in the flesh. After some time of communicating, he began asking for money, claiming various adversities. After he extorted several hundred thousand dollars from her, he disappeared. I find it had to sympathize with her, although the women in the discussion did. First of all, a man should not be asking a woman for money.

Ultimately the fellow contacted the woman, and gave a full confession. Some of the women thought it was wonderful that, at least, the guy confessed and apologized. They can't understand that a psychopath gets a thrill of revealing his scam to his victims.

\paragraph{Prevailing Prayer}


Jesus didn't charge a penny to teach his disciples how to pray while on the mount. Even better, he gave everyone a free lunch afterwards. In our time it might cost you \$297 to learn how to pray. I suspect that the organizers of the Dream Trip Conference\footnote{\url{https://sparks.clickfunnels.com/dream-trip-2019}} are praying that a lot of people sign up.

\paragraph{Correlation and Causation}


Bachelors will often interject ``correlation is not causation'' into a discussion. For the empiricist, aka ``science'', there is nothing but correlation, as David Hume pointed out: causation is purely a mental construction to relate commonly observed phenomena.

Cause therefore is a metaphysical, and not an empirically derived, category.

\paragraph{The Debate over Abortion}


\begin{quotex}
To limit the increase of children, or put to death any of the later progeny is accounted infamous.
\flright{\textsc{Tacitus}, \emph{describing the customs of the Germanic tribes}}
\end{quotex}


The debate over abortion has become quite intense in the USA recently, so it may serve as an example of rhetoric and reputation. In the now infamous interview of Ben Shapiro by the BBC\footnote{\url{https://www.youtube.com/watch?v=6VixqvOcK8E}}, Andrew Neil posited that the Georgia heartbeat abortion law was barbaric and a return to the so-called ``Dark Ages''. Now I don't know if abortion was illegal under King Arthur, but his illegitimate son, Mordred, was not aborted. Nevertheless, the king was perfectly willing to kill Mordred after his birth, an idea that has even been floated by some pro-abortionists. To kill born babies should seem barbaric to everyone, although you should not count on it.

If he meant by the ``Dark Ages'', however, the time when the Roman Empire was taken over by the Germanic tribes, he may be correct. Rome's population had been decreasing due, primarily, to the desire to avoid childbirth. There was discussion about that among the Romans.

On the other hand, the Germanics, as Tacitus noted, encouraged large families. Unlike the Greeks and Romans, who could put to death unwanted or undesired children on orders of the father, those Germanics did not murder their offspring. Readers will have to determine for themselves if the Germanic tribes were more or less moral than the Romans.

Presumably, then, the Germanics did not perform abortions. The issue depends on whether or not the foetus is a human being, at a particular stage of development. This is not an abstract question about when it becomes a ``person'' or is informed by a soul. It is biological reality. Many deny that the foetus is human, claiming instead that it is a ``part'' of the woman, like an unwanted tumour. Those who are more honest recognize the foetus as human but justify abortion because it is ``trespassing'' on the mother against her will. Curiously, the repeated mantra that abortion should be allowed for the ``health of the mother'' implicitly admits that the foetus is really a child and not merely an unwelcome body part. Obviously, a mother requires a child, although no one seems to acknowledge that contradiction.

The other argument of some nun named Joan Chittister is that to be ``pro-life'' means you should be willing to be taxed to support someone else's children. That hardly seems just and I doubt that Sister Joan pays any taxes at all.

The anger among pro-abortionists arises from the change in lifestyle that would necessarily follow. Casual sex would need to be discouraged in a society where abortion was totally outlawed. It is easy to see what other changes to sexual mores would be necessary. Thus the emotional argument for abortion rights are quite strong an persuasive.

Anti-abortion laws as a policy can only work among a virtuous people and will never be accepted in today's USA.


\flrightit{Posted on 2019-05-26 by Cologero}

\begin{center}* * *\end{center}

\begin{footnotesize}
\begin{sffamily}

\texttt{Balrog Booger on 2022-02-04 at 12:55 said:}

I'm always fascinated by the state of the discourse on the abortion question. Almost everyone fails to grasp the deepest implication of this conflict. What women are demanding by full abortion rights is nothing short of sovereignty over the future of the entire human race; de facto matriarchal society. They are demanding that men be completely disenfranchised over the most critical issue of the secular world: what kind of people are allowed to be born.

I wouldn't expect women to be conscious of this, since so much of their behavior is implicit rather than explicit, but men's failure to recognize the stakes is disconcerting. Even in places like the
\_manosphere\_ and \_red-pill\_ groups that specialize in such topics,

the vocabulary and rhetoric they use on the abortion question is still not much evolved beyond the typical political discourse of the normie world. This blindness exists despite the fact that they recognize the aphroditian quality of contemporary society and denounce women's collective quest for power that comes at everyone's expense.

Though, maybe that shouldn't be surprising. The overwhelming majority of people in those Internet spheres live in the simulacrum created by their discourse rather than anything resembling reality. Women have collectively declared war on men. Will men recognize this fact and react accordingly before the catastrophe on the horizon arrives? It seems impossible.

\hfill

\texttt{Cologero on 2022-02-05 at 20:43 said:}

It's a step too far to posit the ``women have collectively declared war on men''. It takes two to tango, as Confucius say, so men participate in the act. Don't forget that pain in childbirth --- understood in a broad sense --- is a punishment, so not all women desire it. As for the other point, women have always had sovereignty over the future of the human race through sexual selection. The alternative is forced marriage and that would be a hard sell.

\hfill

\texttt{Arthur Konrad on 2022-02-06 at 07:01 said:}

@ Cologero

`As for the other point, women have always had sovereignty over the future of the human race through sexual selection.'

That would be akin to saying that Alsace and Lorraine `selected' Germany over France in the aftermath of the Franco-Prussian war. We all know who does the selection. Sovereign is the one who is at liberty to make a choice.

\end{sffamily}
\end{footnotesize}

